\documentclass[10pt,twocolumn,letterpaper]{article}

% CVPR template
\usepackage{cvpr}
\usepackage{times}
\usepackage{epsfig}
\usepackage{graphicx}
\usepackage{amsmath}
\usepackage{amssymb}

%%%%%%%%% PAPER ID
\def\cvprPaperID{****}
\def\confName{COMP 6630 Final Project Proposal}
\def\confYear{2026}

\begin{document}

%%%%%%%%% TITLE
\title{Predicting UFC Fight Outcomes Using Machine Learning}

\author{
Clay Ramey \quad Nicholas Winfrey \quad Muhammad Khattak \\
{\tt\small car0126@auburn.edu, ntw0015@auburn.edu, muk0003@auburn.edu} \\
\\
COMP 5630/6630 -- Machine Learning \\
Auburn University
}

\maketitle
\thispagestyle{empty}

\section{Project Idea}
This project aims to predict UFC fight outcomes using machine learning. Given two fighters and their historical statistics, we will predict the probability that each fighter wins. This is a binary classification problem where the target variable is the fight result (fighter A wins vs.\ fighter B wins), with P(B wins) = 1 $-$ P(A wins).

We will compare multiple ML approaches covered in this course---logistic regression, decision trees, and multi-layer perceptrons (MLPs)---to determine which best captures the factors that influence fight outcomes.

\section{Dataset Description}
Our dataset contains approximately 5,000 UFC fights with detailed fighter information including:

\begin{itemize}
\item \textbf{Fight outcomes}: Winner and method of victory
\item \textbf{Physical attributes}: Height, weight, reach, stance
\item \textbf{Performance history}: Win/loss records, strike statistics, takedown accuracy, submission attempts
\item \textbf{Fight details}: Date, weight class, event information
\end{itemize}

From this raw data, we will engineer features such as win rates, recent form (last 3--5 fights), striking/grappling differentials between fighters, and experience gaps. Preprocessing will include handling missing values, normalizing continuous features, and encoding categorical variables. We will use a temporal train/test split to avoid data leakage.

\section{Proposed ML Task}
\textbf{Input}: Fighter A and Fighter B statistics (individual and relative features).

\textbf{Output}: Predicted probability of each fighter winning.

\textbf{Models}: Logistic regression (baseline), decision trees/random forests, and MLP. Each model will be tuned using cross-validation.

\textbf{Evaluation}: Accuracy, precision, recall, F1-score, confusion matrices, and AIC comparison across models. Feature importance analysis will identify which attributes are most predictive of fight outcomes.

\section{Software}
Python, AWS S3, scikit-learn, matplotlib, pandas, polars, PyTorch, Git.

\section{Work Breakdown}

\textbf{Clay Ramey (Data)}: Data cleaning, preprocessing, and feature engineering---building fighter statistics from raw fight data, handling missing values, and normalization.

\textbf{Muhammad Khattak (Modeling)}: Implementing and training models (logistic regression, decision trees, MLP), hyperparameter tuning and model selection.

\textbf{Nicholas Winfrey (Evaluation \& Integration)}: Model evaluation (accuracy, precision, recall, confusion matrices, AIC comparison), visualizations, and working between the data and modeling sides to refine features and tune parameters based on evaluation results.

\section{Tentative Schedule}
Regular meetings: Mondays at 8:00 PM and Tuesdays at 6:15 PM.

\begin{itemize}
\item \textbf{Weeks 1--2}: Data exploration, cleaning, and feature engineering
\item \textbf{Weeks 3--4}: Baseline model implementation and initial evaluation
\item \textbf{Weeks 5--6}: Model refinement and hyperparameter tuning
\item \textbf{Weeks 7--8}: Final evaluation, analysis, and report writing
\end{itemize}

\end{document}